% Hafta 5 — Yol geçişini kapatan iki adım
% Derleme: py -3.12 tools/latex/derle.py docs/week-5/assets/h05-05-yol-gecisi.tex
\documentclass[tikz,border=8pt]{standalone}
\usepackage{cen429-cizim}
\begin{document}
\begin{tikzpicture}
  \node[baslik] (b) at (0,0) {Yol geçişini kapatan iki adım};
  \node[altbaslik] at (b.south west) {sıra önemli: önce kanonikleştir, SONRA denetle};
  \node[kotukutu=34mm, anchor=north west] (n1) at (0,-2.4)
    {\kb{Ham yol}\\\kod{../../etc/passwd}\\\kod{\%2e\%2e\%2f ...}};
  \node[kutu=36mm, right=9mm of n1] (n2)
    {\kb{1· KANONİKLEŞTİR}\\realpath / normalize\\tek doğru yazım};
  \node[uyarikutu=36mm, right=9mm of n2] (n3)
    {\kb{2· KÖK İÇİNDE Mİ?}\\izinli klasörle\\karşılaştır};
  \node[iyikutu=34mm, right=9mm of n3] (n4)
    {\kb{Kabul / RED}\\kök dışındaysa reddet};
  \draw[kotuok] (n1.east) -- (n2.west);
  \draw[ok] (n2.east) -- (n3.west);
  \draw[iyiok] (n3.east) -- (n4.west);
  \node[fit=(n1)(n2)(n3)(n4), inner sep=0pt] (grp) {};
  \node[dip, text width=155mm, anchor=north west] at ($(grp.south west)+(0,-8mm)$)
    {Ham dize üzerinde yapılan denetim kodlanmış biçimlerle atlatılır.};
\end{tikzpicture}
\end{document}
